\documentclass[11pt]{article}
\usepackage[margin=1in]{geometry}
\usepackage[T1]{fontenc}
\usepackage[utf8]{inputenc}
\usepackage{amsmath,amssymb,amsthm}
\usepackage{enumitem}

\title{ICPC World Finals 2021\\D. Guardians of the Gallery}
\author{}
\date{}

\begin{document}
\maketitle

\section*{Problem Summary}

We are given a simple polygon, the guard position, and the center of a sculpture. The guard may move
only inside the polygon, and we want the minimum distance needed to reach a position from which at
least half of an infinitesimally small circular sculpture is visible.

So the real task is:
\begin{itemize}[leftmargin=*]
    \item describe the region of points from which the sculpture is ``visible enough'',
    \item then find the shortest path from the guard to that region.
\end{itemize}

\section*{Key Observations}

\begin{itemize}[leftmargin=*]
    \item Fix a ray starting at the sculpture center and going outward. As we move along that ray, the
    only moments when the visibility status can change are when the ray hits the polygon boundary.

    \item A \emph{proper crossing} of the ray by the polygon boundary blocks both halves of the sculpture
    at once. A \emph{tangent touch} blocks only one side. Along one ray, the guard may continue moving
    until both sides have been blocked at least once.

    \item Therefore, on every relevant ray, the valid positions form one prefix segment starting at the
    sculpture center and ending at the first point where both sides are blocked.

    \item The only rays that matter are the rays passing through polygon vertices. Between two consecutive
    such rays, the combinatorial description of the blocking boundary does not change.

    \item Once these boundary rays are known, the rest is a standard shortest-path-in-polygon problem:
    a shortest path bends only at polygon vertices, and the final straight segment ends on one of the
    boundary spokes.
\end{itemize}

\section*{Algorithm}

\begin{enumerate}[leftmargin=*]
    \item For every polygon vertex $v$, consider the ray from the sculpture center $s$ through $v$.

    \item For one fixed ray, scan all polygon edges:
    \begin{itemize}[leftmargin=*]
        \item if an edge crosses the ray properly, it blocks both sides;
        \item if a boundary vertex lies on the ray, inspect the two polygon chains incident to that vertex:
        if they approach the ray from opposite sides, this is also a two-sided block; otherwise it is a
        one-sided block.
    \end{itemize}
    Let $t_+$ and $t_-$ be the first distances on the ray where the positive and negative side become
    blocked. The usable part of the ray ends at
    \[
        \max(t_+, t_-).
    \]

    \item Rays with the same direction are merged, keeping only the farthest reachable endpoint. This gives
    a set of boundary spokes, each represented by one segment from $s$ to its cutoff point.

    \item Build the usual visibility graph on the guard and all polygon vertices. Two nodes are connected if
    the straight segment between them lies inside the polygon closure.

    \item Run Dijkstra from the guard to all graph nodes.

    \item For each graph node $q$ and each boundary spoke $T$:
    \begin{itemize}[leftmargin=*]
        \item compute the subset of $T$ directly visible from $q$;
        \item add the distance from $q$ to the nearest visible point of $T$.
    \end{itemize}
    The visible subset of one spoke is found by splitting the spoke at all parameters where the line from
    $q$ to the spoke passes through a polygon vertex. Between two consecutive split points, the visibility
    status is constant, so a midpoint test is enough.

    \item Also check whether the guard itself, or any polygon vertex reached by Dijkstra, already lies in
    the valid region.
\end{enumerate}

\section*{Correctness Proof}

We prove that the algorithm returns the correct answer.

\paragraph{Lemma 1.}
Along a fixed ray starting at the sculpture center, the set of valid guard positions is exactly one prefix
segment of that ray.

\paragraph{Proof.}
Moving outward along the ray can only lose visibility, never regain it, because every blocking event lies
between the sculpture and all farther points on the same ray.

A proper crossing of the boundary blocks both sides of the tiny circle immediately. A tangential contact
blocks only one side. To see at least half of the sculpture, at least one side must remain unblocked.
Hence the valid points are precisely those before the first position where both sides have been blocked.
So the valid set on that ray is one prefix segment. \qed

\paragraph{Lemma 2.}
The endpoint of the valid prefix on a ray is completely determined by polygon edges and vertices lying on
that ray, and only rays through polygon vertices can change the combinatorial answer.

\paragraph{Proof.}
Between two consecutive rays through polygon vertices, every polygon edge stays on the same side of the
ray, so the ordering of one-sided and two-sided blocking events does not change. Therefore the cutoff
value varies continuously without changing combinatorial type, and all changes happen exactly at rays
through vertices. The cutoff on such a ray is determined by the first blocking event on each side, which
comes from the ray hitting polygon edges or polygon vertices. \qed

\paragraph{Lemma 3.}
There exists an optimal path that consists of a shortest path from the guard to some visibility-graph node,
followed by one straight segment to a boundary spoke.

\paragraph{Proof.}
The valid region is contained in the polygon and is bounded by the spokes computed above. Any shortest
path inside a simple polygon bends only at polygon vertices; the last bend before entering the target
region is therefore either the guard itself or a polygon vertex. From that final bend, the remainder of the
path is a straight segment to the first point where the path enters the target region, and that entry point
lies on the boundary of the target region, hence on one of the boundary spokes. \qed

\paragraph{Lemma 4.}
For a fixed graph node $q$ and a fixed spoke $T$, splitting $T$ at parameters where the line $qx$ passes
through a polygon vertex is sufficient to determine all directly visible subsegments of $T$.

\paragraph{Proof.}
As the endpoint $x$ moves continuously along $T$, the segment $qx$ can change its intersection pattern
with the polygon boundary only when it passes through a polygon vertex. Between two consecutive such
events, no combinatorial change is possible, so the visibility status is constant on the whole open
subsegment. Therefore midpoint testing on each interval is sufficient. \qed

\paragraph{Theorem.}
The algorithm outputs the correct answer for every valid input.

\paragraph{Proof.}
By Lemmas 1 and 2, the algorithm computes exactly the boundary spokes of the valid region.
By Lemma 3, some optimal path is represented by a visibility-graph path to a node $q$ and one final
straight segment from $q$ to a spoke. By Lemma 4, the algorithm checks exactly all such final straight
segments. Dijkstra provides the correct shortest distance from the guard to every possible last bend.
Taking the minimum over all graph nodes and all spokes therefore yields the true optimum. \qed

\section*{Complexity Analysis}

Let $n$ be the number of polygon vertices.

\begin{itemize}[leftmargin=*]
    \item Computing all spoke cutoffs costs $O(n^2)$.
    \item Building the visibility graph costs $O(n^3)$ because there are $O(n)$ nodes and each visibility
    test against the polygon costs $O(n)$.
    \item For every graph node and every spoke, we split by all polygon vertices and run $O(n)$ segment
    tests, so this phase is $O(n^4)$.
\end{itemize}

With $n \le 100$, this is easily fast enough. Memory usage is $O(n^2)$.

\section*{Implementation Notes}

\begin{itemize}[leftmargin=*]
    \item Boundary points are allowed. Geometrically, standing exactly on a limiting spoke corresponds to
    seeing exactly half of the infinitesimal circle, which is valid.
    \item The code tests segment containment in the polygon by collecting all boundary-intersection
    parameters, then checking the midpoint of every open interval between them.
    \item Several vertices can lie on the same ray from the sculpture; those rays are merged by normalized
    direction.
\end{itemize}

\end{document}
